% sections/description.tex
% Phase 9: Description du projet et objectifs (RAPP-03)

\section{Description du projet et objectifs}
\label{sec:description}

\subsection{Donn\'ees utilis\'ees}
\label{sec:description_donnees}

Le jeu de donn\'ees utilis\'e dans cette \'etude provient de la table publique
\texttt{crypto\_polygon\allowbreak .transactions} h\'eberg\'ee sur Google
BigQuery~\cite{bigquery_public}. Les transactions ont \'et\'e filtr\'ees pour ne
retenir que les transferts du jeton WETH (Wrapped Ether), identifi\'e par
l'adresse de contrat \texttt{0x7ceB23fD6bC0adD59E62\allowbreak ac25578270cFf1b9f619}~\cite{weth_polygon}, sur la
seule journ\'ee du 5 ao\^ut 2024. Le fichier r\'esultant,
\texttt{polygon\_weth\_crash\_2024-08-05.csv}, contient 401\,709 transactions
pour un volume total de 74\,Mo. Le Tableau~\ref{tab:dataset} r\'esume les
caract\'eristiques principales du jeu de donn\'ees.

\begin{table}[H]
\centering
\caption{Caract\'eristiques du jeu de donn\'ees WETH Polygon.}
\label{tab:dataset}
\begin{tabular}{ll}
\toprule
Propri\'et\'e & Valeur \\
\midrule
Source & Google BigQuery (\texttt{crypto\_polygon}) \\
Token & WETH (Wrapped Ether) \\
Blockchain & Polygon PoS \\
Date & 5 ao\^ut 2024 \\
Transactions & 401\,709 \\
Colonnes & \texttt{transaction\_hash}, \texttt{chrono}, \texttt{source}, \texttt{target}, \texttt{weight} \\
Taille & 74\,Mo \\
\bottomrule
\end{tabular}
\end{table}

Chaque enregistrement du fichier CSV est d\'ecrit par cinq colonnes~:
\texttt{transaction\_hash}, l'identifiant unique de la transaction sur la
blockchain Polygon~; \texttt{chrono}, l'horodatage Unix (nombre de secondes
depuis l'epoch)~; \texttt{source}, l'adresse Ethereum du portefeuille
\'emetteur~; \texttt{target}, l'adresse Ethereum du portefeuille r\'ecepteur~;
et \texttt{weight}, le montant en WETH transf\'er\'e.

\subsection{Questions de recherche}
\label{sec:description_rq}

Le projet s'organise autour de six questions de recherche.

\paragraph{RQ1 --- Analyse de la topologie du r\'eseau.}
Examiner la structure topologique globale du r\'eseau~: distribution des degr\'es
entrants et sortants, densit\'e, composantes faiblement et fortement connexes,
diam\`etre et longueur moyenne des chemins. D\'eterminer si le r\'eseau pr\'esente
des propri\'et\'es de type sans \'echelle, caract\'eristiques des r\'eseaux o\`u
quelques n\oe{}uds concentrent un nombre disproportionn\'e de connexions.

\paragraph{RQ2 --- Analyse de la centralit\'e.}
Identifier les n\oe{}uds les plus influents du r\'eseau par quatre mesures de
centralit\'e~: degr\'e, \textit{betweenness}, \textit{closeness} et PageRank.
Comparer les classements obtenus par chaque mesure et interpr\'eter les r\^oles
des acteurs centraux dans l'\'ecosyst\`eme DeFi (routeurs DEX, ponts
inter-cha\^ines, coffres-forts de protocoles).

\paragraph{RQ3 --- D\'etection de communaut\'es.}
Appliquer l'algorithme de Louvain pour r\'ev\'eler la structure communautaire
sous-jacente du r\'eseau. Mesurer la modularit\'e obtenue et analyser la
distribution des tailles de communaut\'es. \'Evaluer le degr\'e de
compartimentalisation du r\'eseau et ses implications pour la propagation des
chocs financiers.

\paragraph{RQ4 --- Dynamique temporelle.}
D\'ecouper la journ\'ee du 5 ao\^ut 2024 en fen\^etres horaires et suivre
l'\'evolution des m\'etriques du r\'eseau au cours du temps. Identifier les pics
d'activit\'e en nombre de transactions et en volume WETH, puis les relier au
contexte financier mondial (sessions de march\'e asiatique, europ\'eenne et
am\'ericaine).

\paragraph{RQ5 --- Propri\'et\'es petit-monde.}
\'Evaluer le coefficient de clustering et la longueur moyenne des chemins du
r\'eseau. Comparer ces m\'etriques au mod\`ele al\'eatoire d'Erd\H{o}s-R\'enyi
calibr\'e avec les m\^emes param\`etres ($n$ et $p$). Quantifier le caract\`ere
petit-monde par le coefficient $\sigma$ de Humphries et Gurney.

\paragraph{RQ6 --- Analyse des flux pond\'er\'es.}
Analyser la distribution des volumes WETH \'echang\'es sur les ar\^etes du
r\'eseau. Quantifier la concentration des flux par le coefficient de Gini et la
courbe de Lorenz. Identifier les \textit{whales} (portefeuilles r\'ealisant des
transferts de tr\`es grande valeur) et caract\'eriser leur r\^ole dans le r\'eseau.

\subsection{Pr\'etraitement des donn\'ees}
\label{sec:description_preprocessing}

Le pr\'etraitement comprend trois \'etapes. D'abord, les ar\^etes multiples
entre une m\^eme paire d'adresses sont agr\'eg\'ees par somme des poids, avec
un attribut \texttt{tx\_count} conservant le nombre de transactions. Ensuite,
un graphe orient\'e pond\'er\'e (DiGraph) de 12\,880 sommets et 39\,999 ar\^etes
est construit avec NetworkX~; une version non orient\'ee (33\,827 ar\^etes) est
d\'eriv\'ee pour RQ3 et RQ5. Enfin, un filtrage \textit{dust} \`a $10^{-12}$
WETH supprime 1\,598 ar\^etes (4\,\%) de bruit num\'erique pour RQ6.
